\documentclass[12pt]{amsart}


\usepackage{fullpage}
\usepackage{mdframed}
\usepackage{colonequals}
\usepackage{algpseudocode}
\usepackage{algorithm}
\usepackage{tcolorbox}
\usepackage[all]{xy}
\usepackage{proof}
\usepackage{mathtools}
\usepackage{bbm}
\usepackage{amssymb}
\usepackage{amsthm}
\usepackage{amsmath}
\usepackage{amsxtra}
\usepackage{enumitem}
\usepackage{euler}



% Theorem environments
\newtheorem{theorem}{Theorem}[section]
\newtheorem{theorem*}{Theorem}
\newtheorem{definition}[theorem]{Definition}
\newtheorem{corollary}{Corollary}[theorem]
\newtheorem{lemma}[theorem]{Lemma}
\newtheorem{proposition}[theorem]{Proposition}
\newtheorem*{remark}{Remark}
\newtheorem*{claim}{Claim}
\newtheorem*{example}{Example}
\newtheorem*{warning}{Warning}

\newtheorem*{exercisehelper}{Exercise.}
\newenvironment{exercise}[1]{%
  \IfBlankTF{#1}
    {\renewcommand{\exercisehelper}{\textbf{Exercise} \unskip}}
    {\renewcommand\exercisehelper{\textbf{Exercise #1}}}%
  \exercisehelper
}{\endexercisehelper}

\theoremstyle{remark}
\newtheorem*{solution}{Solution}
\newcommand{\mathcat}[1]{\textup{\textbf{\textsf{#1}}}} % for defined terms

\newenvironment{problem}[1]
{\begin{tcolorbox}\noindent\textbf{Problem #1}.}
{\vskip 6pt \end{tcolorbox}}

\newenvironment{enumalph}
{\begin{enumerate}\renewcommand{\labelenumi}{\textnormal{(\alph{enumi})}}}
{\end{enumerate}}

\newenvironment{enumroman}
{\begin{enumerate}\renewcommand{\labelenumi}{\textnormal{(\roman{enumi})}}}
{\end{enumerate}}

\let\H\relax
\let\P\relax
\let\Pr\relax
\let\d\relax
\let\limsup\relax
\let\1\relax
\let\div\relax


\DeclarePairedDelimiter\ceil{\lceil}{\rceil}
\DeclarePairedDelimiter\floor{\lfloor}{ \rfloor}
\newcommand{\fr}[2]{\left(\frac{#1}{#2}\right)}
\newcommand{\norm}[1]{\left\vert\left\vert#1\right\vert\right\vert}
\newcommand{\dbar}{\overline{\partial}}
\newcommand{\d}{\partial}
\newcommand{\RP}{\mathbb{RP}}
\newcommand{\CP}{\mathbb{CP}}
\newcommand{\tbigwedge}{{\textstyle{\bigwedge}}}


\newcommand{\bbni}{\bigbreak \noindent}
\newcommand{\todo}[1]{\textcolor{red}{\textbf{TODO:} #1}}
\newcommand{\red}[1]{\textcolor{red}{#1}}


% Blackboard Font
\newcommand{\N}{\mathbb N}
\newcommand{\Z}{\mathbb Z}
\newcommand{\Q}{\mathbb Q}
\newcommand{\R}{\mathbb R}
\newcommand{\C}{\mathbb C}
\newcommand{\F}{\mathbb F}
\newcommand{\G}{\mathbb G}
\newcommand{\H}{\mathbb H}
\newcommand{\bb}{\mathbb}

% Calligraphic font
\newcommand{\CC}{\mathcal{C}}
\newcommand{\DD}{\mathcal{D}}
\newcommand{\EE}{\mathcal{E}}
\newcommand{\FF}{\mathcal{F}}
\newcommand{\GG}{\mathcal{G}}
\newcommand{\HH}{\mathcal{H}}
\newcommand{\LL}{\mathcal{L}}
\newcommand{\OO}{\mathcal{O}}
\newcommand{\II}{\mathcal{I}}


% Linear Algebra
\DeclareMathOperator{\GL}{GL}
\DeclareMathOperator{\PGL}{PGL}
\DeclareMathOperator{\Mat}{Mat}
\DeclareMathOperator{\Id}{Id}
\DeclareMathOperator{\Tr}{Tr}
\DeclareMathOperator{\tr}{tr}
\DeclareMathOperator{\Nm}{Nm}
\DeclareMathOperator{\opspan}{span}
\DeclareMathOperator{\rk}{rk}
\DeclareMathOperator{\sgn}{sgn}
\DeclareMathOperator{\Sym}{Sym}
\DeclareMathOperator{\Alt}{Alt}
\DeclareMathOperator{\codim}{codim}



% Category Theory / Homological Algebra
\DeclareMathOperator{\id}{id}
\DeclareMathOperator{\colim}{colim}
\DeclareMathOperator{\Aut}{Aut}
\DeclareMathOperator{\End}{End}
\DeclareMathOperator{\Hom}{Hom}
\DeclareMathOperator{\Mor}{Mor}
\DeclareMathOperator{\Ob}{Ob}
\DeclareMathOperator{\Ab}{Ab}
\DeclareMathOperator{\Coh}{Coh}
\DeclareMathOperator{\QCoh}{QCoh}
\DeclareMathOperator{\ann}{ann}
\DeclareMathOperator{\img}{img}
\DeclareMathOperator{\coker}{coker}
\DeclareMathOperator{\Ext}{Ext}
\DeclareMathOperator{\Tor}{Tor}


% Algebraic Geometry
\newcommand{\A}{\mathbb A}
\newcommand{\P}{\mathbb P}
\newcommand{\V}{\mathbb V}
\newcommand{\I}{\mathbb I}
\DeclareMathOperator{\div}{div}
\DeclareMathOperator{\Spec}{Spec}
\DeclareMathOperator{\Proj}{Proj}
\DeclareMathOperator{\Frob}{Frob}
\DeclareMathOperator{\Pic}{Pic}
\DeclareMathOperator{\Weil}{Weil}


% Unclassfified
\DeclareMathOperator{\Cinf}{C^{\infty}}
\DeclareMathOperator{\Ell}{Ell}
\DeclareMathOperator{\CL}{\mathcal{CL}}
\DeclareMathOperator{\Log}{Log}
\DeclareMathOperator{\Arg}{Arg}
\DeclareMathOperator{\Res}{Res}
\DeclareMathOperator{\val}{val}

\DeclareMathOperator{\Supp}{Supp}
\DeclareMathOperator{\cl}{cl}
\DeclareMathOperator{\disc}{disc}
\DeclareMathOperator{\M}{M}
\DeclareMathOperator{\opchar}{char}
\DeclareMathOperator{\argmax}{argmax}
\DeclareMathOperator{\argmin}{argmin}
\DeclareMathOperator{\limsup}{limsup}
\DeclareMathOperator{\Ind}{Ind}
\DeclareMathOperator{\Pr}{\mathbf{Pr}}
\DeclareMathOperator{\Exp}{\mathbb{E}}
\DeclareMathOperator{\Var}{\mathbf{Var}}

\setlength{\parindent}{0pt}
\setlength{\parskip}{5pt}
\setlength{\hfuzz}{4pt}




\begin{document}

\title{Degrees of Hodge Loci}

\author{Talk by David Urbanik}
\maketitle

\section{Setup}
Let $U_d$ be the algebraic variety parametrizing smooth degree $d$ hypersurfaces in $\P^3$. This is $\P(\Sym^d(k^\vee))$. 

\begin{theorem}[Lefschetz]
    There exists a countable collection $\{z_j\}_{j \in J}$ of closed subvarieties of $U_d$ such that, 
    \[ \operatorname{NL} = \bigcup_{j \in J} z_j\]
    This is the Noether-Lefschetz. 
\end{theorem}

The NL locus is the collection of hypersurfaces $H_b$ which contain an irreducible curve $C \subset H_b$, not arising as a complete intersection in $\P^3$ with a surface of this type. 

We have a question: \red{Can we understand the degrees of the components $z_j$}. 

If $d = 4$, there is a polynomial, 
\begin{enumerate}
    \item $Q_1$ such that $H_b$ contains a $\P^1$ if and only if $Q_1(b) = 0$ and $\deg Q_1 = 320$. 
    \item $Q_2$ such that $H_b$ contains a quadric and $\deg Q_2 = 2508$. 
\end{enumerate}

To do this, we do an interlude into Hodge Theory. 

\section{Hodge Theory}
We will reformulate this using Hodge Structures on $H^2(H_b, \Z)$. 

\begin{definition}
    Given an integral lattice $V$, a (polarized) Hodge structure on $V$ of weight $k$ is a direct sum decomposition: 
    \[ V_\C = \bigoplus_{p+q = k} H^{p,q}\]
    satisfying $H^{p,q} = \overline{H^{q,p}}$. 
\end{definition}

The idea is that these come from a direct sum decomposition of de Rham cohomology with complex coefficients on a complex manifold. 

If $X$ is a complex manifold, locally on $X$ there aresmooth complexified differential forms that break into holomorphic and anti-holomorphic components. For a Kahler manifold, even the global ones, the ones that give cohomology, descend with this decomposition with canonical representatives.

\begin{conjecture}[Hodge Conjecture.]
    If $X$ is a smooth projective algebraic variety then the cohomology classes in $H^{2p}(X, \Q) \cap H^{p, p}(X)$ are in the $\Q$-span of algebraic cycle classes.
\end{conjecture}

\begin{align*}
    H^2(H_b, \Z) \to H^2(H_b, \Z)_{prim} \oplus \Z
\end{align*}
with $H_b \subset \P^3$, $C \subset H_b$, and $[C]_{prim} \in H^2(X_b, \Z)$. 

We define,
\[ Q : H^2(H_b, \Z)_{prim} \oplus H^2(H_b, \Z)_{prim} \to \Z\]     

\begin{definition}
    We define
    \[ \operatorname{NL}_q = \{ b \in \operatorname{NL} : \exists C \subset H_b, Q([c]_{prim}, [c]_{prim}) = q\} \]
\end{definition}


\begin{theorem}[Urbanik.]
    There exists $C > 0$ depeneding on $d$ such that, 
    \[ \deg \operatorname{NL}_q \leq c q^e\]
    with 
    \begin{align*}
        e = (b_2-1)^2 + 2h^{2,0}((b_2-1) - h^{2,0}) - 2
    \end{align*}
    where $b_2 = \dim_\C H^2(H_b, \Q)$ and $h^{2,0} = \dim_\C H^{2,0}$.
\end{theorem}

Next, we talk about the method. 

\section{Seigel Sets}
A problem about degrees is a problem about counting points. If we took a sufficiently general locus complimentary codim $F$ and intersected it with $NL_q$, we get the degree by counting intersections. 

The idea is: We reduce the question to producing upper bounds for loci where there are extra Hodge vectors. 

For Seigel Sets, we consider a space $D$ which is the parameter space for all (polarized) Hodge structures on $V \cong \Z^{2n}$ with fixed Hodge numbers $h^{2,0}$, $h^{1,1}$, and $h^{0,2}$. 

There is a natural discrete group $\Gamma$ such that, 
\[ D/\Gamma = \{ \text{pol. Hodge structures on V} \}/\text{iso}\]

\begin{definition}
    A Siegel set is a kind of fundamental set for the covering
    \[ \pi: D \to D/\Gamma\]
    it is associated to a triple $(P, S, K)$ where: 
    \begin{enumerate}
        \item $P$ is a minimal $\Q$-parabolic subgroup(?)
        \item $S$ is the maximal $\Q$-split torus in $P$, stable under the Cartan involutions associated to $K$.
        \item $K$ is a maximal compact subgroup corresponding to a choice $h \in D$. 
    \end{enumerate}
\end{definition}

For $G = \GL_n$, $K = O(n)$, $S$ is the diagonal matrices, and $P$ is the lower triangular matrices. 

A Siegel set in $G(\R)$ is a product, 
\begin{align*}
    \Omega A_t K \subset G(\R)
\end{align*}
where $\Omega$ is compact in $P(\R)$. 

If one constructs Siegel sets for $G = \Aut(V, Q)$ associated to $D$, 
\begin{align*}
    G(\R) &\to D \qquad g \mapsto g\cdot h \\
    \Omega A_t K &\mapsto T
\end{align*}
where $T$ is a Hodge-theoretic Siegel set. 

The upshort is: One can choose Hodge theoretic Siegel sets such that, 
\begin{enumerate}
    \item They give a fundamental set for $\pi: D \to D/\Gamma$. 
    \item The set of $h \in T$ that admit an integeral Hodge vector $v$ with $Q(v,v) \leq q$ has size that is is $\leq$ some polynomial in $q$. 
\end{enumerate}


\end{document}